\documentclass[12pt, a4paper]{article}

% 引入必要的宏包
\usepackage[UTF8]{ctex}     % 中文支持
\usepackage{amsmath, amssymb, amsthm} % 数学符号和环境
\usepackage{geometry}       % 页面设置
\usepackage{enumitem}       % 列表定制
\usepackage{graphicx}       % 图片支持
\usepackage{fancyhdr}       % 页眉页脚

% 页面边距设置
\geometry{left=2.5cm, right=2.5cm, top=2.5cm, bottom=2.5cm}

% 宏定义：方便排版选择题选项
% 使用方法：\choices{选项A}{选项B}{选项C}{选项D}
\newcommand{\choices}[4]{
    \begin{enumerate}[label=\Alph*., itemsep=0pt, parsep=0pt, topsep=5pt, labelsep=0.5em]
        \item #1
        \item #2
        \item #3
        \item #4
    \end{enumerate}
}
% 横向排列的选项 (2列)
\newcommand{\choicesTwo}[4]{
    \begin{enumerate}[label=\Alph*., itemsep=0pt, parsep=0pt, topsep=5pt, labelsep=0.5em]
        \begin{minipage}{0.45\linewidth} \item #1 \end{minipage}
        \begin{minipage}{0.45\linewidth} \item #2 \end{minipage}
        \begin{minipage}{0.45\linewidth} \item #3 \end{minipage}
        \begin{minipage}{0.45\linewidth} \item #4 \end{minipage}
    \end{enumerate}
}
% 横向排列的选项 (4列)
\newcommand{\choicesFour}[4]{
    \begin{enumerate}[label=\Alph*., itemsep=0pt, parsep=0pt, topsep=5pt, labelsep=0.5em]
        \begin{minipage}{0.22\linewidth} \item #1 \end{minipage}
        \begin{minipage}{0.22\linewidth} \item #2 \end{minipage}
        \begin{minipage}{0.22\linewidth} \item #3 \end{minipage}
        \begin{minipage}{0.22\linewidth} \item #4 \end{minipage}
    \end{enumerate}
}

\title{\textbf{《高等数学》必刷习题——基础版}}
\author{}
\date{}

\begin{document}

\section*{第一章 函数（基础）}

\begin{enumerate}
    \item 以下区间是函数 $ y = \sin x $ 的单调递增区间的是 \underline{\hspace{1cm}}。
    \choicesFour{$ [0,\frac{\pi}{2}] $}{$ [0,\pi] $}{$ [\frac{\pi}{2},\pi] $}{$ [\pi,\frac{3\pi}{2}] $}

    \item $ y = \frac{e^{x} - e^{-x}}{2} $ 是 \underline{\hspace{1cm}}。
    \choicesFour{奇函数}{偶函数}{非奇非偶函数}{无法确定}

    \item 函数 $ y = |x \cos(-x)| $ 是 \underline{\hspace{1cm}}。
    \choicesFour{有界函数}{偶函数}{单调函数}{周期函数}

    \item 设 $ f(x) $ 的定义域为 $ [1,2] $ ，则函数 $ f(x^{2}) $ 的定义域是 \underline{\hspace{1cm}}。
    \choicesFour{$ [1,2] $}{$ [1,\sqrt{2}] $}{$ \left[-\sqrt{2},\sqrt{2}\right] $}{$ \left[-\sqrt{2},-1\right] \cup \left[1,\sqrt{2}\right] $}

    \item 下列各对函数中，\underline{\hspace{1cm}} 中的两个函数是相同的。
    \choices
    {$ f(x)=\frac{x^{2}-1}{x-1} $, $ g(x)=x+1 $}
    {$ f(x)=\sqrt{x^{2}} $, $ g(x)=x $}
    {$ f(x) = \ln x^{2} $, $ g(x) = 2\ln x $}
    {$ f(x) = \sin^{2}x + \cos^{2}x $, $ g(x) = 1 $}

    \item 若 $ f(x) = \log_{a} x \quad (a > 0, a \neq 1) $ ，则 $ f(x) + f(y) $ 等于 \underline{\hspace{1cm}}。
    \choicesFour{$ f(x+y) $}{$ f(xy) $}{$ f(x-y) $}{$ f(\frac{y}{x}) $}

    \item 下列给定区间中是函数 $ f(x)=|x^{2}-1| $ 的单调有界区间的是 \underline{\hspace{1cm}}。
    \choicesFour{$ [-1,1] $}{$ (1,+\infty) $}{$ [-2,-1] $}{$ [-2,0] $}

    \item 不等式 $ \frac{|x-1|-1}{|x-3|}>0 $ 的解集(用区间表示)为 \underline{\hspace{1cm}}。
    \choicesTwo
    {$ (- \infty, 0) $}
    {$ (- \infty, 3) \cup (3, +\infty) $}
    {$ (2,3) \cup (3, +\infty) $}
    {$ (- \infty, 0) \cup (2, 3) \cup (3, +\infty) $}

    \item 已知函数 $ f(x) $ 的定义域为 $ [0,1] $ ，则函数 $ f(\ln x) $ 的定义域是 \underline{\hspace{1cm}}。
    \choicesFour{$ [1,e] $}{$ [e,+\infty) $}{$ (0,e] $}{$ (1,e) $}

    \item 假设函数 $ f(x) = \sin\frac{x}{2} + \cos\frac{x}{3} $ ，则 $ f(x) $ 的周期为 \underline{\hspace{2cm}}。

    \item 如果函数 $ f(x) $ 的定义域是 $ \left[-\frac{1}{3},3\right] $ ，则 $ f(\frac{1}{x}) $ 的定义域是 \underline{\hspace{2cm}}。

    \item 函数 $ f(x) = x \frac{a^{x}-1}{a^{x}+1} $ 的图像关于 \underline{\hspace{2cm}} 对称。

    \item 求 $ y=\frac{2x}{\sqrt{2-x}}+\ln(1+x) $ 的定义域 \underline{\hspace{2cm}}。

    \item 设 $ f(x) $ 是奇函数， $ g(x) $ 是偶函数，则 $ f(x)g(x) $ 是 \underline{\hspace{2cm}}。

    \item 函数 $ y=\sqrt{x^{2}-x-6}-\arcsin\frac{2x-3}{9} $ 的定义域是 \underline{\hspace{2cm}}。
\end{enumerate}


\end{document}